% sim_model_rsi.sv ve rsi_pkg.sv - Ayrıntılı kod açıklaması (taslak)
% Overleaf: pdfLaTeX. Satır numaraları v1.0.9 kaynak dosyalarına göredir.
% "TODO" işaretli yerler doldurulacak.
\documentclass[11pt,a4paper]{article}

\usepackage[T1]{fontenc}
\usepackage[utf8]{inputenc}
\usepackage[turkish]{babel}
\usepackage{lmodern}
\usepackage[margin=2.5cm]{geometry}
\usepackage{booktabs}
\usepackage{longtable}
\usepackage{array}
\usepackage{enumitem}
\usepackage{xcolor}
\usepackage{amssymb}
\usepackage{textcomp}
\usepackage{tikz}
\usetikzlibrary{arrows.meta,positioning}
\usepackage{listings}
\usepackage[hidelinks]{hyperref}
\usepackage{fancyhdr}

\setlength{\parindent}{0pt}
\setlength{\parskip}{0.6em}
\renewcommand{\arraystretch}{1.2}
\pagestyle{fancy}
\fancyhf{}
\lhead{RSİ modeli}
\rhead{Kod açıklaması (taslak)}
\cfoot{\thepage}

\newcommand{\us}{\,\textmu{}s}
\newcommand{\kod}[1]{\texttt{#1}}
\newcommand{\todo}[1]{\textcolor{red!60!black}{\textbf{[TODO]} #1}}
\newcommand{\satir}[1]{\textcolor{black!60}{\scriptsize(satır #1)}}

\lstset{
  basicstyle=\ttfamily\footnotesize,
  breaklines=true,
  frame=single,
  columns=fullflexible,
  keepspaces=true,
  upquote=true,
  extendedchars=false
}

\tikzset{
  kutu/.style={draw=black!60, rounded corners=3pt, align=center,
               minimum height=1cm, inner sep=4pt, font=\small},
  ana/.style={kutu, draw=blue!70!black, fill=blue!8, very thick},
  ok/.style={-{Stealth[length=2.2mm]}, thick, black!70},
  etiket/.style={font=\scriptsize, text=black!70}
}

\title{\textbf{RSİ Modeli}\\[0.3em]\large \kod{sim\_model\_rsi.sv} ve \kod{rsi\_pkg.sv}
       kod açıklaması (taslak)}
\author{}
\date{\today}

\begin{document}
\maketitle
\tableofcontents
\newpage

%---------------------------------------------------------------
\section{Genel bakış}

\todo{Modelin rolü: gerçek RSİ'nin yerine geçen doğrulayıcı. SİFPGA'nın PCIe RQ
üzerinden yazdığı bellek yazımlarını (MemWr) toplar, MSI1/MSI2 gelince "turu kapatır":
belleği çözer, doğrular, dosyaya döker ve belleği temizler.}

\begin{itemize}[nosep]
  \item Dosyalar: \kod{src/sim\_model\_rsi.sv} (607 satır), \kod{src/rsi\_pkg.sv}
        (199 satır), \kod{src/cpu\_pkg.sv} içindeki \kod{rq\_assembler} ve
        \kod{rq\_pkt} sınıfları.
  \item \kod{RSI\_ID} (1..3): her RSİ için bir model örneği. Çıktı:
        \kod{<OUT\_PREFIX>\_rsi<RSI\_ID>.txt} ve ortak \kod{rsi\_tur\_ozeti.txt}.
  \item Portlar \kod{sim\_model\_cpu} ile aynı; CQ tarafı boştur, HB üretilmez.
\end{itemize}

\subsection{Veri akışı şeması}

\begin{figure}[h]
\centering
\begin{tikzpicture}[node distance=0.9cm and 1.0cm]
  \node[kutu, text width=2.6cm] (sif) {SİFPGA\\ PCIe RQ};
  \node[kutu, text width=2.8cm, right=of sif] (asm) {\kod{rq\_assembler}\\ beat $\to$ TLP};
  \node[kutu, text width=2.6cm, right=of asm] (mem) {\kod{mem}\\ seyrek kelime belleği};
  \node[ana, text width=3.0cm, below=1.4cm of mem] (rnd) {\kod{process\_round}\\ çöz, doğrula};
  \node[kutu, text width=2.8cm, left=of rnd] (isr) {\kod{run\_isr}\\ MSI1..MSI4};
  \node[kutu, text width=2.6cm, left=of isr] (msi) {MSI girişi\\ \kod{cfg\_interrupt\_msi\_int}};
  \node[kutu, text width=3.0cm, below=1.1cm of rnd] (pkg) {\kod{rsi\_pkg}\\ tur durumu, özet};
  \draw[ok] (sif) -- (asm);
  \draw[ok] (asm) -- (mem);
  \draw[ok] (msi) -- (isr);
  \draw[ok] (isr) -- (rnd);
  \draw[ok] (mem) -- (rnd);
  \draw[ok] (rnd) -- (pkg);
\end{tikzpicture}
\caption{Modelin iç veri akışı. \todo{Şemayı kodla karşılaştırıp düzelt.}}
\end{figure}

%---------------------------------------------------------------
\section{Parametreler ve portlar}
\satir{27--62}

\begin{longtable}{@{}p{3.6cm}p{2.2cm}p{9.1cm}@{}}
\toprule
Parametre & Varsayılan & Anlamı \\
\midrule
\endhead
\kod{RSI\_ID} & 1 & Bu örneğin RSİ numarası (1..3) \\
\kod{OUT\_PREFIX} & \kod{rsi\_out} & Çıktı dosyası ön eki (\kod{+RSI\_OUT=} ile değişir) \\
\kod{ISR\_LATENCY\_NS} & 100 & MSI'dan sonra ISR'nin başlaması için bekleme (ns) \\
\kod{MAX\_WORDS} & 4194304 & Bellekteki en çok 4 baytlık kelime; aşılırsa yazma atılır \\
\kod{DETAIL\_LINES} & 64 & Dosyaya yazılan en fazla örnek satırı (0: yok, $-1$: hepsi) \\
\kod{FP\_SCALE\_TOL} & 0 & Ayrılmış; float doğruluğu tam olmak zorunda \\
\kod{CQ\_KEEP\_WIDTH}, \kod{RQ\_KEEP\_WIDTH} & 8 & AXI-S keep genişliği \\
\bottomrule
\end{longtable}

\todo{Portları grupla: saat/reset, CQ (boş), RQ (girdi), MSI, HB. \kod{s\_axis\_rq\_tready}
her zaman \kod{4'hF}; \kod{cfg\_interrupt\_msi\_fail} her zaman 0.}

%---------------------------------------------------------------
\section{RQ yakalama ve bellek}
\satir{83--141}

\subsection{\kod{rq\_assembler}}
\todo{Her kabul edilen beat'i (tvalid ve tready) toplar; tlast'ta 16 baytlık
descriptor çözülür. DW0: adres[31:2], DW1: adres[63:32], DW2: DWORD sayısı ve istek
tipi. tuser: ilk beat'te first\_be, last\_be.}

\subsection{\kod{store\_pkt}}
\todo{Payload'ı 4 baytlık küçük-endian kelimelere çevirip \kod{mem[(addr>>2)+i]}'ye
yazar. Anahtar: kelime adresi. \kod{partial\_be\_warn}, \kod{dropped} sayaçları.
\kod{MemWr} dışındaki tipler uyarıyla atılır.}

\begin{lstlisting}[language={}]
// kod parçası buraya (satır 126-141)
\end{lstlisting}

%---------------------------------------------------------------
\section{MSI mekanizması ve ISR}
\satir{143--195}

\todo{Adımlar: (1) \kod{msi\_int} yükselen kenar algılama, (2) 2 clk gecikmeli
\kod{cfg\_interrupt\_msi\_sent}, (3) \kod{msi\_pend} birikimi ve \kod{irq\_vec}'e
aktarım, (4) ISR iş parçacığı: \kod{irq\_flag} bekle, \kod{asm.collecting} bitene
kadar bekle, \kod{ISR\_LATENCY\_NS} bekle, vektörü al, \kod{run\_isr}.}

\begin{table}[h]
\centering
\begin{tabular}{@{}ll@{}}
\toprule
Vektör biti & İşlem \\
\midrule
bit0 (MSI1) & Tur kapanır, \kod{basarili = 1} \\
bit1 (MSI2) & Tur kapanır, \kod{basarili = 0} \\
bit2 (MSI3) & \kod{msi3\_pending} artar, loga yazılır \\
bit3 (MSI4) & Yalnızca loga yazılır \\
bit4..31 & Tanımsız bit uyarısı \\
\bottomrule
\end{tabular}
\caption{\kod{run\_isr} davranışı. \todo{Doğrula.}}
\end{table}

%---------------------------------------------------------------
\section{Yardımcı fonksiyonlar}
\satir{197--290}

\subsection{float32 $\to$ tam sayı (\kod{f32\_ok}, \kod{f32\_val})}
\todo{Neden tam olmayan değer reddediliyor? Üs, mantis ve kesir biti maskesi.}

\subsection{Desen kodu (\kod{code\_ok}, \kod{code\_m}, \kod{code\_n})}
\todo{$\mathrm{kod}=((m+1)\ll15)\,|\,(n\,\&\,\mathtt{0x7FFF})$; geçerli aralık
$1\le m+1\le 960$.}

\subsection{\kod{wd}, \kod{count\_valid}, \kod{ks}, \kod{wfmt}}
\todo{\kod{wd(i)}: mesajın i. kelimesi (yoksa 0). \kod{count\_valid(k)}: çarpan k
için geçerli (kanal, örnek) çifti sayısı. \kod{ks}: denenen çarpanlar (3, 1, 2).}

\subsection{Mesaj düzeni sabitleri}
\begin{table}[h]
\centering
\begin{tabular}{@{}lll@{}}
\toprule
Sabit & Değer (kelime) & Anlamı \\
\midrule
\kod{HDR\_WORDS} & 24 & ID + LEN + HEADER 88~B \\
\kod{GAC\_WORDS} & 480 & GAC SCIT 1920~B \\
\kod{SHOB\_WORDS} & 8 & SHOB SCIT 32~B \\
\kod{SIB\_WORDS} & 4 & SİBFPGA SCIT 16~B \\
\kod{FOOT\_WORDS} & 4 & Alt bilgi 16~B \\
\kod{TAIL\_WORDS} & 496 & Kuyruk toplamı \\
\bottomrule
\end{tabular}
\end{table}

%---------------------------------------------------------------
\section{\kod{process\_round}: turu kapatma}
\satir{294--592}

\todo{Bu bölüm kodun kalbi. Her alt bölüm: ne yapıyor, hangi değişkenleri kullanıyor,
hangi log satırını yazıyor, hata/uyarı sayacını ne zaman artırıyor.}

\subsection{Hazırlık ve anahtarlar}
\satir{309--324} \todo{Yerel durum temizliği, \kod{mem.first/next} ile artan adres
sırası, \kod{nwords}, \kod{nbytes}, \kod{prev\_nbytes}, \kod{msi3\_pending}.}

\subsection{Hüzme ID'siyle tur eşleştirme (v1.0.9)}
\satir{326--347} \todo{\kod{wd(2)} hüzme ID'si; \kod{t\_hid[j]} ve \kod{t\_nk}
ile ilk uyan tur seçilir; aradakiler \kod{rsi\_report\_missing}. Neden sıralı
eşleştirme yetmiyordu?}

\subsection{Başlangıç adresi ve bitişiklik}
\satir{366--389} \todo{\kod{start\_expected = rsi\_start\_addr + t\_off}; bitişiklik
kontrolü.}

\subsection{Ham döküm}
\satir{391--394} \todo{İlk 32 kelime her zaman yazılır.}

\subsection{ID ve LEN kontrolü, teşhis}
\satir{396--420} \todo{\kod{lendiff}, \kod{TESHIS} satırları, beklenen toplam bayt.}

\subsection{HEADER çözümü}
\satir{422--464} \todo{Kelime düzeni tablosu: durum, alma hüzme sayısı, darbe,
örnek, GAC maskeleri, SHOB kartları, ayrılmış alanlar; \kod{1DEA0003} ile karşılaştırma.}

\subsection{Kanal sayısı $K$}
\satir{466--482} \todo{$K=(\text{bayt}-2080)/(8N)$ ve beklenenle karşılaştırma.}

\subsection{Birleştirme kuralı seçimi}
\satir{484--503} \todo{Çarpan 3, 1, 2 denenir; ilk tam uyan seçilir.}

\subsection{Kanal kanal çözme}
\satir{505--546} \todo{Dizin: \kod{HDR\_WORDS + c*2*N + n} (I), \kod{+N} (Q). Kontroller:
$I>0$, $I \bmod k=0$, $Q=-I$, kod geçerli, kanal içinde aynı $m$, $n$ sırası.}

\subsection{Hüzme listesi}
\satir{548--562}

\subsection{SCIT bölümleri ve alt bilgi}
\satir{564--584}

\subsection{Sonuç ve rapor}
\satir{586--592} \todo{\kod{dogru = (errors==0)}; \kod{rsi\_report\_round}; \kod{reset\_round}.}

%---------------------------------------------------------------
\section{\kod{rsi\_pkg}: ortak durum}
\satir{rsi\_pkg.sv}

\subsection{Veri yapıları}
\todo{\kod{NRSI}=3, \kod{MAXR}=256, \kod{MAXB}=1024; düz diziler (xsim kısıtı).
\kod{t\_*} turlar, \kod{r\_*} raporlar.}

\subsection{\kod{track\_cfg\_msgs}}
\todo{CPU sürücüsünün her TLP için çağırdığı izleyici; \kod{1DEA0001} (RSİ başlangıç
adresleri) ve \kod{1DEA0003} (tur parametreleri) bayt konumları.}

\subsection{\kod{rsi\_report\_round}}
\todo{Tüm RSİ'ler raporlayınca özet satırı, birleşim kontrolü (her hüzme 0..M-1 bir
kez), VERİ GELMEDİ ve LENFARK/MSI3 notları.}

\subsection{\kod{rsi\_report\_missing}, \kod{rsi\_flush\_missing}, yardımcılar}

%---------------------------------------------------------------
\section{Çıktı dosyaları}

\todo{\kod{rsi\_out\_rsiN.txt} satır türleri (ISR, TUR, ESLESME, TESHIS, ADRES OK,
HAM DOKUM, MESAJ, HEADER, KANAL SAYISI, BIRLESTIRME KURALI, HUZME SAYISI, GAC SCIT,
FOOTER, SONUC) ve \kod{rsi\_tur\_ozeti.txt} biçimi. Örnek log parçası ekle.}

%---------------------------------------------------------------
\section{Sınırlar ve bilinen davranışlar}

\begin{itemize}[nosep]
  \item Bellek her turda \kod{reset\_round} ile temizlenir; MSI'siz yazımlar bir
        sonraki tura sızabilir.
  \item En çok 256 tur (\kod{MAXR}) izlenir.
  \item Float değerler tam sayı olmak zorundadır.
  \item MSI3, MSI4 yalnızca kaydedilir; hata enjeksiyonu yok.
  \item \todo{Diğer.}
\end{itemize}

\end{document}
